\subsection{\texorpdfstring{THE COEFFICIENTS \(\mathbf{N}_{\alpha \beta}\)}{THE COEFFICIENTS \textbackslash mathbf\{N\}\_\{\textbackslash alpha \textbackslash beta\}}}

In this number, we again denote by \((\mathfrak{g},\mathfrak{h})\) a split semi-simple Lie algebra.

Lemma 2. There exists a family \((X_{\alpha})_{\alpha \in \mathbb{R}}\) such that, for all \(\alpha \in \mathbb{R}\) ,

\[
X _ {\alpha} \in \mathfrak {g} ^ {\alpha} \text { and } [ X _ {\alpha}, X _ {- \alpha} ] = - H _ {\alpha}.
\]

Let \(R_{1}\) be a subset of R such that \(R = R_{1} \cup (-R_{1})\) and \(R_{1} \cap (-R_{1}) = \varnothing\) . For \(\alpha \in R_{1}\) , choose an arbitrary non-zero element \(X_{\alpha}\) of \(g^{\alpha}\) . There exists a unique \(X_{-\alpha} \in g^{-\alpha}\) such that \([X_{\alpha}, X_{-\alpha}] = -H_{\alpha}\) (Th. 1 (iv)). Then

\[
\left[ X _ {- \alpha}, X _ {\alpha} \right] = H _ {\alpha} = - H _ {- \alpha}.\tag{Q.E.D.}
\]

If \((X_{\alpha})_{\alpha \in \mathbb{R}}\) is one family satisfying the conditions of Lemma 2, the most general family satisfying these conditions is \((t_{\alpha}X_{\alpha})_{\alpha \in \mathbb{R}}\) where \(t_{\alpha}\in k^{*}\) and \(t_{\alpha}t_{-\alpha} = 1\) for all \(\alpha \in \mathbb{R}\) .

In the remainder of this number, we denote by \((X_{\alpha})_{\alpha\in\mathbb{R}}\) a family satisfying the conditions of Lemma 2. We denote by \(\langle\cdot,\cdot\rangle\) a non-degenerate invariant symmetric bilinear form on g.

Every \(x \in \mathfrak{g}\) can be written uniquely in the form

\[
x = h + \sum_ {\alpha \in \mathrm{R}} \mu_ {\alpha} X _ {\alpha} \quad (h \in \mathfrak {h}, \mu_ {\alpha} \in k).
\]

The bracket of two such elements can be calculated by means of the following formulas:

\[
[ h, X _ {\alpha} ] = \alpha (h) X _ {\alpha}
\]

\[
[ X _ {\alpha}, X _ {\beta} ] = \left\{ \begin{array}{l l} 0 & \text {if} \alpha + \beta \notin \mathrm{R} \cup \{0 \} \\ - H _ {\alpha} & \text {if} \alpha + \beta = 0 \\ \mathrm{N} _ {\alpha \beta} X _ {\alpha + \beta} & \text {if} \alpha + \beta \in \mathrm{R} \end{array} \right.
\]

the \(N_{\alpha\beta}\) being non-zero elements of k.

Lemma 3. For all \(\alpha \in \mathbb{R}\) ,

\[
\langle X _ {\alpha}, X _ {- \alpha} \rangle = - \frac {1}{2} \langle H _ {\alpha}, H _ {\alpha} \rangle .
\]

Indeed,

\[
\begin{array}{c} 2 \langle X _ {\alpha}, X _ {- \alpha} \rangle = \langle \alpha (H _ {\alpha}) X _ {\alpha}, X _ {- \alpha} \rangle = \langle [ H _ {\alpha}, X _ {\alpha} ], X _ {- \alpha} \rangle \\ = \langle H _ {\alpha}, [ X _ {\alpha}, X _ {- \alpha} ] \rangle = - \langle H _ {\alpha}, H _ {\alpha} \rangle . \end{array}
\]

Lemma 4. Let \(\alpha, \beta \in \mathrm{R}\) be such that \(\alpha + \beta \in \mathrm{R}\) . Let \(p\) (resp. \(q\) ) be the largest integer \(j\) such that \(\beta + j\alpha \in \mathrm{R}\) (resp. \(\beta - j\alpha \in \mathrm{R}\) ). Then,

\[
\mathrm{N} _ {\alpha , \beta} \mathrm{N} _ {- \alpha , \alpha + \beta} = - p (q + 1)\tag{3}
\]

\[
\mathrm{N} _ {- \alpha , \alpha + \beta} \langle H _ {\beta}, H _ {\beta} \rangle = - \mathrm{N} _ {- \alpha , - \beta} \langle H _ {\alpha + \beta}, H _ {\alpha + \beta} \rangle\tag{4}
\]

\[
\mathrm{N} _ {\alpha , \beta} \mathrm{N} _ {- \alpha , - \beta} = (q + 1) ^ {2}.\tag{5}
\]

Let \(\rho\) be the representation of \(\mathfrak{sl}(2,k)\) on \(\mathfrak{g}\) defined by \(X_{\alpha}\) . The element \(e = X_{\beta + p\alpha}\) is primitive of weight \(p + q\) (Prop. 4 (i)). Put

\[
e _ {n} = \frac {(- 1) ^ {n}}{n !} \rho (X _ {-}) ^ {n} e \quad \text { for } n \geq 0.
\]

By Prop. 1 of §1,

\[
(\mathrm{ad} X _ {\alpha}) e _ {p} = (q + 1) e _ {p - 1}
\]

\[
(\operatorname{ad} X _ {- \alpha}) (\operatorname{ad} X _ {\alpha}) e _ {p} = - p (q + 1) e _ {p}.
\]

This proves (3) since \(e_p\) is a non-zero element of \(\mathfrak{g}^{\beta}\) .

The form \(\langle \cdot, \cdot \rangle\) being invariant, we have

\[
\langle [ X _ {- \alpha}, X _ {\alpha + \beta} ], X _ {- \beta} \rangle = - \langle X _ {\alpha + \beta}, [ X _ {- \alpha}, X _ {- \beta} ] \rangle
\]

SO

\[
\mathrm{N} _ {- \alpha , \alpha + \beta} \langle X _ {\beta}, X _ {- \beta} \rangle = - \mathrm{N} _ {- \alpha , - \beta} \langle X _ {\alpha + \beta}, X _ {- \alpha - \beta} \rangle
\]

which, in view of Lemma 3, proves (4).

The restriction of \(\langle\cdot,\cdot\rangle\) to h is non-degenerate and invariant under the Weyl group (Prop. 5). Identify h and \(h^{*}\) by means of this restriction. If \(\gamma\in R\) , \(H_{\gamma}\) is identified with \(2\gamma/\langle\gamma,\gamma\rangle\) (Chap. VI, §1, no. 1, Lemma 2); hence, for all \(\gamma,\delta\in R\) ,

\[
\frac {\langle \gamma , \gamma \rangle}{\langle \delta , \delta \rangle} = \frac {\langle H _ {\delta} , H _ {\delta} \rangle}{\langle H _ {\gamma} , H _ {\gamma} \rangle}.\tag{6}
\]

Now, by Chap. VI, §1, no. 3, Prop. 10,

\[
\frac {\langle \alpha + \beta , \alpha + \beta \rangle}{\langle \beta , \beta \rangle} = \frac {q + 1}{p}\tag{7}
\]

so, by (3), (4), (6), (7),

\[
\begin{array}{r} \mathrm{N} _ {\alpha , \beta} \mathrm{N} _ {- \alpha , - \beta} = - \mathrm{N} _ {\alpha , \beta} \mathrm{N} _ {- \alpha , \alpha + \beta} \frac {\langle H _ {\beta} , H _ {\beta} \rangle}{\langle H _ {\alpha + \beta} , H _ {\alpha + \beta} \rangle} \\ = - \mathrm{N} _ {\alpha , \beta} \mathrm{N} _ {- \alpha , \alpha + \beta} \frac {q + 1}{p} = (q + 1) ^ {2}. \end{array}
\]

DEFINITION 3. A Chevalley system for \((\mathfrak{g},\mathfrak{h})\) is a family \((X_{\alpha})_{\alpha \in \mathbb{R}}\) such that

\begin{enumerate}
\def\labelenumi{(\roman{enumi})}
\item
  \(X_{\alpha} \in g^{\alpha}\) for all \(\alpha \in R;\)
\item
  \([X_{\alpha}, X_{-\alpha}] = -H_{\alpha}\) for all \(\alpha \in \mathrm{R}\) ;
\item
  the linear map from \(\mathfrak{g}\) to \(\mathfrak{g}\) which is equal to \(-1\) on \(\mathfrak{h}\) and which takes \(X_{\alpha}\) to \(X_{-\alpha}\) for all \(\alpha \in \mathrm{R}\) is an automorphism of \(\mathfrak{g}\) .
\end{enumerate}

The extension of this definition to the case where \((\mathfrak{g},\mathfrak{h})\) is split reductive is immediate.

We shall show (§4, no. 4, Cor. of Prop. 5) that Chevalley systems for \((\mathfrak{g},\mathfrak{h})\) exist.

PROPOSITION 7. Let \((X_{\alpha})_{\alpha \in \mathbb{R}}\) be a Chevalley system for \((\mathfrak{g},\mathfrak{h})\) . We retain the notation of Lemma 4. Then, \(\mathrm{N}_{-\alpha, - \beta} = \mathrm{N}_{\alpha ,\beta}\) and \(\mathrm{N}_{\alpha ,\beta} = \pm (q + 1)\) for \(\alpha ,\beta ,\alpha +\beta \in \mathbb{R}\) .

Let \(\varphi\) be the automorphism of \(\mathfrak{g}\) considered in Def. 3 (iii). Then

\[
\begin{array}{r} \mathrm{N} _ {- \alpha , - \beta} X _ {- \alpha - \beta} = [ X _ {- \alpha}, X _ {- \beta} ] = [ \varphi (X _ {\alpha}), \varphi (X _ {\beta}) ] = \varphi ([ X _ {\alpha}, X _ {\beta} ]) \\ = \varphi (\mathrm{N} _ {\alpha , \beta} X _ {\alpha + \beta}) = \mathrm{N} _ {\alpha , \beta} X _ {- \alpha - \beta} \end{array}
\]

so \(N_{-\alpha,-\beta}=N_{\alpha,\beta}\) . Now \(N_{\alpha,\beta}=\pm(q+1)\) by (5).

PROPOSITION 8. Let \((X_{\alpha})_{\alpha \in \mathbb{R}}\) be a Chevalley system for \((\mathfrak{g},\mathfrak{h})\) . Let \(\mathrm{M}\) be a \(\mathbf{Z}\) -submodule of \(\mathfrak{h}\) containing the \(H_{\alpha}\) and contained in the group of weights of \(R^{\vee}\) . Let \(\mathfrak{g}_{\mathbf{Z}}\) be the Z-submodule of g generated by M and the \(X_{\alpha}\) . Then \(\mathfrak{g}_{\mathbf{Z}}\) is a Z-Lie subalgebra of g, and the canonical map from \(g_{Z} \otimes_{Z} k\) to g is an isomorphism.

If \(\alpha, \beta \in \mathbb{R}\) are such that \(\alpha + \beta \in \mathbb{R}\) , then \(\mathrm{N}_{\alpha, \beta} \in \mathbf{Z}\) (Prop. 7). On the other hand, if \(\alpha \in \mathbb{R}\) and \(h \in \mathbb{M}\) , then \(\alpha(h) \in \mathbf{Z}\) (Chap. VI, §1, no. 9). This proves that \(\mathfrak{g}_{\mathbf{Z}}\) is a \(\mathbf{Z}\) -Lie subalgebra of \(\mathfrak{g}\) . On the other hand, M is a free abelian group of rank dim \(\mathfrak{h}\) (Algebra, Chap. VII, §3, Th. 1), so \(\mathfrak{g}_{\mathbf{Z}}\) is a free abelian group of rank dim \(\mathfrak{g}\) ; this implies the last assertion.

